% =====================================================================
%  ANEXOS
%  Cada \chapter{...} después de \appendix se numera como Anexo A, B, C...
% =====================================================================
\clearpage
\capituloSinNumero{Anexos}

Esta sección reúne el material de apoyo que complementa el trabajo. Cada anexo
se identifica con una letra y se cita en el texto por su nombre.

\appendix
\renewcommand{\appendixname}{Anexo}
\renewcommand{\chaptername}{Anexo}

\titleformat{\chapter}[display]
  {\normalfont\bfseries\Large}{Anexo \thechapter}{10pt}{\LARGE}
\addtocontents{toc}{\protect\addvspace{6pt}}

\chapter{Instrumentos de recolección de datos}
\label{anx:instrumentos}

Incluya aquí los cuestionarios, guías de entrevista, fichas de observación u
otros instrumentos utilizados.

\chapter{Material complementario}
\label{anx:complementario}

Incluya tablas extensas, transcripciones, gráficos adicionales, normas o
cualquier documento que apoye el trabajo pero cuya presencia en el cuerpo
interrumpiría la lectura. Cada anexo debe mencionarse al menos una vez en el
texto, por ejemplo: «véase el Anexo~\ref{anx:complementario}».
